\chapter*{Abstract}
\addcontentsline{toc}{chapter}{Abstract}

A mobile telecom operator keeps six years of billing records in a single Oracle table. Monthly
and audit reports slow down as the table grows, purging old data is a long \code{DELETE} that
produces gigabytes of redo, and the legal retention rules are applied by hand. As two
consultants, we were asked to design and evaluate a data lifecycle strategy that combines
partitioning, indexing, compression and archiving, and to prove its effect with measurements.

We built the solution on Oracle Database 19c Enterprise Edition (in Docker) and compared it
with the client's current design on the same synthetic dataset of 23.76 million billing rows
over 72 months. The optimised design partitions the billing table by month with interval
partitioning. A PL/SQL package (\code{DLM\_PKG}) moves every month through four tiers. Hot
months stay uncompressed and writable. Warm months are compressed and read-only. Archived
months are moved to an archive table by two partition exchanges, without copying a row. Months
past the 60-month legal retention are dropped. Each operation is guarded, logged, resumable
after a crash, and validated against control totals captured when the data was loaded.

Measured results: the month-end revenue report reads 54 times fewer blocks and runs 14 times
faster; an audit query on an archived month is 43 times faster. Archiving one month takes
0.42~s instead of 26~s, and purging one takes 0.02~s instead of 5.1~s. The full monthly job,
including the compression step, is 8.9 times faster and produces 3.8 times less redo.
Retained data uses 44\,\% less storage. All 62 retained months match their control totals.
The solution passes 43 automated functional tests and 25 security checks.

\paragraph{Keywords:} Oracle Database 19c, interval partitioning, partition exchange, data
lifecycle management, archiving, basic table compression, partition pruning, retention policy,
PL/SQL, performance benchmarking.
